% answers_short.tex —— 2007 年·普通高等学校招生考试·湖北卷·理科 · 简短答案速查
% 由 tools/make_short_answers.py 自 answers.tex 的 \gksolhead/\gkanswer 提取，
% 并与 2007/2007_数学真题_湖北理_答案要点.md 逐题比对后生成，勿手改（改源 MD 后重跑）。

\begin{gknotice}
\textbf{使用说明}：本节只给最终答案便于快速核对，完整推导见解析版。答案为编写者独立推导，\textbf{非官方标准答案}。
\end{gknotice}

\gkpart{选择题}
\begin{center}\begin{tabular}{ccccccccccc}
\toprule
题号 & 1 & 2 & 3 & 4 & 5 & 6 & 7 & 8 & 9 & 10 \\
\midrule
答案 & B & A & B & D & C & B & A & D & A & A \\
\bottomrule
\end{tabular}\end{center}

\gkpart{填空题}
\begin{itemize}[leftmargin=2.6em]
  \item[\textbf{11.}] $a=6$；$b=\dfrac12$。
  \item[\textbf{12.}] $(2,1)$（答案不唯一）。
  \item[\textbf{13.}] $-\dfrac32$。
  \item[\textbf{14.}] $\dfrac{15}{128}$。
  \item[\textbf{15.}] （1）$y=\begin{cases}10t,&0\leqslant t\leqslant0.1,\\ \left(\dfrac1{16}\right)^{t-0.1},&t>0.1;\end{cases}$（2）$0.6$。
\end{itemize}

\gkpart{解答题}
\begin{itemize}[leftmargin=2.6em]
  \item[\textbf{16.}] （1）$\left[\dfrac{\pi}4,\dfrac{\pi}2\right]$；（2）最大值 $3$，最小值 $2$。
  \item[\textbf{17.}] （1）频率依次为 $0.04$，$0.25$，$0.30$，$0.29$，$0.10$，$0.02$；（2）$0.69$，$0.44$；（3）$1.41$。
  \item[\textbf{18.}] （1）证明见解析；（2）$\left(0,\dfrac{\pi}4\right)$。
  \item[\textbf{19.}] （1）$2\sqrt2p^2$；（2）存在 $l:y=\dfrac p2$，弦长恒为 $p$。
  \item[\textbf{20.}] （1）$b=\dfrac52a^2-3a^2\ln a$，最大值 $\dfrac32\mathrm e^{2/3}$；（2）证明见解析。
  \item[\textbf{21.}] （1）（2）证明见解析；（3）$n=2$ 或 $n=3$。
\end{itemize}
